\par\Needspace{6\baselineskip}\medskip
\hypertarget{native-114-article085}{}\pdfbookmark[2]{85.}{article-085}
\begin{center}85.\end{center}
We shall retain the symbols $r$, $r'$, $N$, $N'$, $p$, $e$, $\Pi$ with the same meanings as above. We denote the difference $N'-N$ by $\Delta$ and the time in which the celestial body moves from the earlier position to the later one by $t$. It is clear that if an approximate value of any one of the quantities $p$, $e$, $\Pi$ is known, the other two can be determined from it. The methods explained in the first section then give the time corresponding to the motion from the first position to the second. If this equals the prescribed time $t$, the assumed value of $p$, $e$, or $\Pi$ is the true one, and the orbit itself has already been found. Otherwise, repeating the calculation with another value differing slightly from the first will show how much the time changes for a small change in the value of $p$, $e$, or $\Pi$. Simple interpolation then gives a corrected value. If the calculation is repeated with that value, the resulting time will either agree exactly with the prescribed time or differ very little from it. With further corrections, we can therefore certainly attain as close an agreement as the logarithmic and trigonometric tables permit.

The problem is thus reduced to showing how, when the orbit is still entirely unknown, at least an approximate value of one of the quantities $p$, $e$, $\Pi$ can be \SourceLeaf{115}determined. We shall now give a method that yields the value of $p$ so precisely that for small values of $\Delta$ no further correction is needed. The first calculation then determines the entire orbit as precisely as the ordinary tables permit. We shall hardly ever need to use this method except for moderate values of $\Delta$. Given the intricacy of the problem, determining a still entirely unknown orbit can scarcely be attempted except with observations that are fairly close together or, more precisely, with observations between which the heliocentric motion is fairly small.
